\section*{The Propers: Notable and Quotable}
{\small
\begin{itemize}
\item \textbf{“Why do you spend your money for that which is not bread?”\enspace\elementref{First Reading}} Henry David Thoreau's \emph{Walden}, “Economy,” invokes Isaiah's question while attacking lives consumed by unnecessary expenditure. A prophetic summons to covenant nourishment becomes an ironic instrument of American economic criticism. The verbal dependence is explicit; Thoreau's program is not treated as Isaiah's own.
\item \textbf{“Neither height nor depth.”\enspace\elementref{Second Reading}} The paired extremes have passed into memorial and devotional inscriptions as a compressed assertion that no spatial or cosmic reach exceeds divine love. The later use turns Paul's apocalyptic catalogue into a portable language of consolation; it does not remove the tribulation named in the passage.
\item \textbf{Five loaves and two fish.\enspace\elementref{Gospel}} The numerical image has become a recurrent name and emblem for food ministries. The narrative's inadequate meal is redirected into an institutional claim that small donated resources, coordinated in service, may feed many. Such names are reception, not evidence that each ministry reproduces the miracle.
\end{itemize}}

\section*{The Propers: Interpretive Possibilities}
\noindent\emph{These exploratory proposals were developed in an AI-assisted editorial process. They are attributed to none of the cited authorities and claim no historical compositional intent.}

\subsection*{The disciples' inventory can become an examination of scarcity speech\enspace\elementref{Gospel; First Reading}}
The disciples' accurate count and Isaiah's invitation without price create a tension between necessary accounting and final judgment. The mechanism is not contempt for budgets but refusal to let an inventory decide whom compassion may receive. Read separately, Isaiah can become abstraction and Matthew mere wonder; together they ask whether scarcity language describes constraints or prematurely dismisses persons. The fruit is prudent generosity. No checked source established this exact conjunction; planning remains necessary, and no institution may promise miraculous multiplication.

\subsection*{The open hand moves through many hands\enspace\elementref{Responsorial Psalm; Gospel; Prayer over Offerings}}
Psalm 145 attributes satisfaction to God's hand; Matthew places blessed bread in disciples' hands; the offerings prayer asks that worshippers become an enduring gift. The proposed mechanism is mediated generosity: divine agency does not erase creaturely service but makes it possible. It reveals distribution as part of worship rather than a sequel detached from it. Chrysostom supplies a near analogue in the disciples' training. The limit is decisive: human agents must not claim divine ownership, coerce gifts, or hide failures behind providential language.

\subsection*{The two Communion branches disclose different answers to hunger\enspace\elementref{Communion A; Communion B; Gospel}}
Wisdom's branch looks backward to manna interpreted as heavenly pedagogy; John's branch centers desire on Christ himself. Compared without conflation, they offer historical remembrance and personal encounter as distinct conclusions to Matthew's meal. The fruit is a richer account of eucharistic hunger. The branches remain mutually exclusive in enactment, John 6 is not Matthew 14, and no claim is made that one antiphon was selected in 2026.

\subsection*{Inseparability can free charity from outcome control\enspace\elementref{Second Reading; Gospel; Prayer after Communion}}
The feeding ends in visible abundance, but Romans locates security in inseparable love rather than guaranteed outcomes; the final prayer still asks continued protection. Together they may free service from the demand that every faithful act visibly succeed. The fruit is perseverance without triumphalism. No exact precedent was located in the checked corpus. This proposal cannot excuse incompetence, preventable scarcity, withheld aid, or refusal to evaluate results.
